\addcontentsline{toc}{section}{Список использованных источников}
\begin{thebibliography}{99}

\bibitem{popovlist}
Попов С.Г. Лист самостоятельной проверки хода выполнения курсовой
работы по дисциплине «Теоретические основы баз данных». 4 семестр.

\bibitem{postgresql16}
PostgreSQL 16 Documentation [Электронный ресурс]: The PostgreSQL
Global Development Group.
URL: \url{https://www.postgresql.org/docs/16} (Дата обращения: 09.05.2026).

\bibitem{takahashi}
Takahashi M., Azuma S., Trend C.L. The Manga Guide to Databases. — No
Starch Press, 2009. — 214 p.

\bibitem{pgx}
pgx - PostgreSQL Driver and Toolkit
URL: \url{https://github.com/jackc/pgx}
(Дата обращения: 20.04.2026).

\bibitem{axenta}
Axenta. Мониторинг транспорта в каршеринге и корпоративном автопарке. — 2026. URL: \url{https://axenta.tech/blog/novosti/monitoring-transporta-v-karsheringe-i-korporativnom-avtoparke/} (Дата обращения: 20.03.2026).

\bibitem{navixy}
Navixy. Что такое телематика? Подробный гид для IoT, GPS и транспорта. — 2020. URL: \url{https://www.navixy.com/ru/blog/what-is-telematics/} (Дата обращения: 20.03.2026).

\bibitem{rbc}
РБК. Совместная мобильность: чего ждать от рынка каршеринга России. — 2024. — URL: \url{https://www.rbc.ru/industries/news/661e8a3a9a79478be5eab43d} (Дата обращения: 20.03.2026).

\bibitem{wiki}
Система управления базами данных [Электронный ресурс]: Википедия. — URL: \url{https://ru.wikipedia.org/wiki/%D0%A1%D0%B8%D1%81%D1%82%D0%B5%D0%BC%D0%B0_%D1%83%D0%BF%D1%80%D0%B0%D0%B2%D0%BB%D0%B5%D0%BD%D0%B8%D1%8F_%D0%B1%D0%B0%D0%B7%D0%B0%D0%BC%D0%B8_%D0%B4%D0%B0%D0%BD%D0%BD%D1%8B%D1%85}
(Дата обращения: 20.03.2026).
\end{thebibliography}